\newcommand{\titulus}{\nomenFesti{S. Monicæ.}
\dies{Die 27. Augusti.}}
\newcommand{\oratio}{\pars{Oratio.}

\noindent Deus, mæréntium consolátor, qui beátæ Mónicæ pias lácrimas in conversióne fílii sui Augustíni misericórditer suscepísti, da nobis, utriúsque intervéntu, peccáta nostra deploráre et grátiæ tuæ indulgéntiam inveníre.

\pars{Pro pace in universo mundo.} \scriptura{Sir. 50, 25; 2 Esdr. 4, 20; \textbf{H416}}

\vspace{-4mm}

\antiphona{II D}{temporalia/ant-dapacemdomine.gtex}

\vfill

\noindent Deus, a quo sancta desidéria, recta consília et iusta sunt ópera: da servis tuis illam, quam mundus dare non potest, pacem; ut et corda nostra mandátis tuis dédita, et hóstium subláta formídine, témpora sint tua protectióne tranquílla.

\noindent Per Dóminum nostrum Iesum Christum, Fílium tuum, qui tecum vivit et regnat in unitáte Spíritus Sancti, Deus, per ómnia sǽcula sæculórum.

\noindent \Rbardot{} Amen.}
\newcommand{\invitatorium}{\pars{Invitatorium.}

\vspace{-4mm}

\antiphona{IV}{temporalia/inv-mirabileminsanctis.gtex}}
\newcommand{\hymnusmatutinum}{\pars{Hymnus.}

\cuminitiali{II}{temporalia/hym-HaecFemina.gtex}}
\newcommand{\lectioi}{\pars{Lectio I.} \scriptura{Ier. 3, 1-5}

\noindent De libro Ieremíæ prophétæ.

\noindent Factum est verbum Dómini ad me, dicens:

\noindent «Si dimíserit vir uxórem suam, et recédens ab eo dúxerit virum álterum, numquid revertétur ad eam ultra? Numquid non pollúta et contamináta est terra illa?

\noindent Tu autem fornicáta es cum amatóribus multis, et revérteris ad me?, dicit Dóminus. Leva óculos tuos ad colles et vide, ubi non prostráta sis. In viis sedébas exspéctans eos quasi Arabs in solitúdine; et polluísti terram in fornicatiónibus tuis et in malítia tua.

\noindent Quam ob rem prohíbitæ sunt stillæ pluviárum, et serótinus imber non fuit. Frons mulíeris meretrícis facta est tibi; noluísti erubéscere.

\noindent Nonne ámodo vocas me: “Pater meus, dux adulescéntiæ meæ tu es! Numquid irascétur in perpétuum aut perseverábit in finem?”. Ecce locúta es et fecísti mala et prævaluísti.}
\newcommand{\responsoriumi}{\pars{Responsorium 1.} \scriptura{\Rbardot{} Iob 20, 27 \Vbardot{} Ps. 68, 23; \textbf{H182}}

\vspace{-5mm}

\responsorium{I}{temporalia/resp-revelabuntcaeli-CROCHU.gtex}{}}
\newcommand{\lectioii}{\pars{Lectio II.} \scriptura{Ier. 3, 19-25; 4, 1-4}

\noindent Ego autem dixi: Quómodo ponam te in fíliis et tríbuam tibi terram desiderábilem, hereditátem præclaríssimam inter gentes?

\noindent Et dixi: Patrem vocábitis me et post me íngredi non cessábitis.

\noindent Sed, quómodo contémnit múlier amatórem suum, sic contempsístis me, domus Israel», dicit Dóminus.

\noindent Vox in cóllibus audíta est, plorátus et supplicátio filiórum Israel, quóniam iníquam fecérunt viam suam, oblíti sunt Dómini Dei sui.

\noindent «Convertímini, fílii, qui avérsi estis a me, et sanábo aversiónes vestras».

\noindent «Ecce nos venímus ad te; tu enim es Dóminus Deus noster. Vere mendáces erant colles et tumúltus móntium; vere in Dómino Deo nostro salus Israel.

\noindent Confúsio comédit labórem patrum nostrórum ab adulescéntia nostra, greges eórum et arménta eórum, fílios eórum et fílias eórum.

\noindent Dormiémus in confusióne nostra, et opériet nos ignomínia nostra, quóniam Dómino Deo nostro peccávimus nos et patres nostri ab adulescéntia nostra usque ad hanc diem et non audívimus vocem Dómini Dei nostri».

\noindent «Si convérteris, Israel, ait Dóminus, ad me convértere; si abstúleris abominatiónes tuas a fácie mea, non effúgies.

\noindent Et iurábis: “Vivit Dóminus!” in veritáte et in iudício et in iustítia, et benedicéntur in ipso gentes et in ipso gloriabúntur.

\noindent Hæc enim dicit Dóminus viro Iudæ et Ierúsalem: Nováte vobis novále et nolíte sérere super spinas. Circumcidímini Dómino et auférte præpútia córdium vestrórum, viri Iudæ et habitatóres Ierúsalem, ne forte egrediátur ut ignis indignátio mea et succendátur, et non sit qui exstínguat, propter malítiam óperum vestrórum».}
\newcommand{\responsoriumii}{\pars{Responsorium 2.} \scriptura{\Rbardot{} Michææ 6, 8 \Vbardot{} Ps. 36, 6; \textbf{H418}}

\vspace{-5mm}

\responsorium{V}{temporalia/resp-indicabotibihomo-CROCHU.gtex}{}}
\newcommand{\lectioiii}{\pars{Lectio III.} \scriptura{Lib. 6, 1: CCL 27,73}

\noindent Ex Confessiónum libris sancti Augustíni epíscopi.

\noindent Iam vénerat ad me mater pietáte fortis, terra maríque me sequens et in perículis ómnibus de te secúra. Nam et per marína discrímina ipsos nautas consolabátur, a quibus rudes abýssi viatóres, cum perturbántur, consolári solent: póllicens eis perventiónem cum salúte, quia hoc ei tu per visum pollícitus eras. Et invénit me, periclitántem quidem gráviter desperatióne indagándæ veritátis.

\noindent Sed tamen ei cum indicássem non me quidem iam esse manichǽum, sed neque cathólicum christiánum, non quasi inopinátum áliquid audíerit, exilívit lætítia; cum iam secúra fíeret ex ea parte misériæ meæ in qua me tamquam mórtuum sed resuscitándum tibi flebat, et féretro cogitatiónis offerébat, ut díceres fílio víduæ: \emph{iúvenis, tibi dico, surge,} et revivésceret et incíperet loqui et tráderes illum matri suæ.

\noindent Nulla ergo turbulénta exultatióne trepidávit cor eius, cum audísset ex tanta parte iam factum quod tibi cotídie plangébat ut fíeret, veritátem me nondum adéptum sed falsitáti iam eréptum; immo vero quia certa erat et quod restábat te datúrum, qui totum promíseras, placidíssime et péctore pleno fidúciæ respóndit mihi crédere se in Christo, quod priúsquam de hac vita emigráret, me visúra esset fidélem cathólicum. Et hoc quidem mihi. Tibi autem, fons misericordiárum, preces et lácrimas densióres, ut acceleráres adiutórium tuum et illumináres ténebras meas, et studiósius ad ecclésiam cúrrere et in Ambrósii ora suspéndi, ad fontem \emph{saliéntis aquæ in vitam ætérnam.}}
\newcommand{\responsoriumiii}{\pars{Responsorium 3.} \scriptura{\Rbardot{} Ps. 44, 5 \Vbardot{} ibid., 8; \textbf{H383}}

\vspace{-5mm}

\responsorium{VIII}{temporalia/resp-propterveritatemetmansuetudinem-CROCHU-cumdox.gtex}{}

\rubrica{vel ad libitum:}

\vspace{3mm}

\pars{Responsorium 3.} \scriptura{Is. 61, 10; \textbf{H113}}

\vspace{-5mm}

\responsorium{VI}{temporalia/resp-induitmedominus-CROCHU-cumdox.gtex}{}}
\newcommand{\hymnuslaudes}{\pars{Hymnus}

\antiphona{II}{temporalia/hym-NobilemChristi.gtex}}
\newcommand{\lectiobrevis}{\pars{Lectio Brevis.} \scriptura{Rom. 12, 1-2}

\noindent Obsecro vos, fratres, per misericórdiam Dei, ut exhibeátis córpora vestra hóstiam vivéntem, sanctam, Deo placéntem, rationábile obséquium vestrum; et nolíte conformári huic sǽculo, sed transformámini renovatióne mentis, ut probétis quid sit volúntas Dei, quid bonum et bene placens et perféctum.}
\newcommand{\responsoriumbreve}{\pars{Responsorium breve.} \scriptura{Ps. 45, 6}

\cuminitiali{VI}{temporalia/resp-adiuvabiteam.gtex}}
\newcommand{\benedictus}{\pars{Canticum Zachariæ.} \scriptura{Ap. 21, 4}

\vspace{-4mm}

\antiphona{I f}{temporalia/ant-abstergetdeus.gtex}

\vspace{-2mm}

\scriptura{Lc. 1, 68-79}

\vspace{-2mm}

\cantusSineNeumas
\initiumpsalmi{temporalia/benedictus-initium-i-f-auto.gtex}

%\vspace{-1.5mm}

\input{temporalia/benedictus-i-f.tex} \Abardot{}}
\newcommand{\benedicamuslaudes}{\cuminitiali{}{temporalia/benedicamus-memoria-laudes.gtex}}
\include{hebdomadaxxi}
\include{feriav}
